\documentclass[aspectratio=169]{beamer}
\usepackage[utf8]{inputenc}
\usepackage[portuguese]{babel}
\usepackage{graphicx}
\usepackage{hyperref}
\usepackage{xcolor}
\usepackage{tikz}
\usepackage{multicol}

% Configuração do tema
\usetheme{Madrid}
\usecolortheme{default}

% Configurações de cores
\definecolor{azul}{RGB}{0,51,102}
\definecolor{azulclaro}{RGB}{51,102,153}
\definecolor{verde}{RGB}{34,139,34}
\definecolor{laranja}{RGB}{255,140,0}
\definecolor{roxo}{RGB}{138,43,226}
\definecolor{vermelho}{RGB}{220,20,60}
\definecolor{amarelo}{RGB}{255,215,0}
\setbeamercolor{title}{fg=white,bg=azul}
\setbeamercolor{frametitle}{fg=white,bg=azulclaro}

% Informações do documento
\title[ChatGPT para Estudar]{Minicurso: Usando o ChatGPT para Estudar}
\subtitle{Dicas para Estudar Melhor com Inteligência Artificial}
\author{Professor Dr. Bento Rafael Siqueira}
\institute{UFLA - Universidade Federal de Lavras\\Departamento de Ciência da Computação}
\date{\today}

\begin{document}

% Slide de título
\begin{frame}
    \titlepage
\end{frame}

\section{Introdução}

\begin{frame}
    \frametitle{Quem gosta de fofoca?}
    \begin{center}
        \Large \textbf{O ChatGPT é como aquela tia fofoqueira!}
    \end{center}
    
    \vspace{0.6cm}
    \begin{center}
    \begin{tikzpicture}[scale=1.4]
        % Cabeça
        \filldraw[fill=orange!30, draw=orange, very thick] (0,0) circle (1.8cm);
        % Olhos
        \filldraw[fill=black] (-0.6,0.4) circle (0.25cm);
        \filldraw[fill=black] (0.6,0.4) circle (0.25cm);
        % Boca sorrindo
        \draw[very thick, red] (-0.7,-0.4) arc (180:0:0.7cm);
        % Cabelo
        \filldraw[fill=brown!50, draw=brown] (-1.4,1.4) .. controls (-1.8,2.2) and (1.8,2.2) .. (1.4,1.4);
    \end{tikzpicture}
    \end{center}
    
    \vspace{0.4cm}
    \begin{center}
    \begin{tikzpicture}[scale=0.85]
        % Caixa 1
        \filldraw[fill=azul!25, draw=azul, very thick, rounded corners=12pt] (-4.5,0.5) rectangle (-1,2.2);
        \node[align=center] at (-2.75,1.6) {\normalsize \textbf{Sabe de tudo!}\\
        \small (ou quase tudo)\\
        \tiny Como a tia que conhece};
        \node[align=center] at (-2.75,0.8) {\tiny todo mundo};
        
        % Caixa 2
        \filldraw[fill=verde!25, draw=verde, very thick, rounded corners=12pt] (-0.5,0.5) rectangle (3,2.2);
        \node[align=center] at (1.25,1.6) {\normalsize \textbf{Adora conversar!}\\
        \small Responde qualquer coisa\\
        \tiny Nunca fica sem assunto};
        \node[align=center] at (1.25,0.8) {\tiny como a tia fofoqueira};
        
        % Caixa 3
        \filldraw[fill=laranja!25, draw=laranja, very thick, rounded corners=12pt] (3.5,0.5) rectangle (7,2.2);
        \node[align=center] at (5.25,1.6) {\normalsize \textbf{24 horas!}\\
        \small Nunca dorme\\
        \tiny Disponível sempre};
        \node[align=center] at (5.25,0.8) {\tiny como a tia no WhatsApp};
        
        % Caixa 4 - Mas cuidado
        \filldraw[fill=amarelo!30, draw=orange, very thick, rounded corners=12pt] (-4.5,-2.5) rectangle (7,-0.2);
        \node[align=center] at (1.25,-1.6) {\normalsize \textbf{Mas cuidado:}};
        \node[align=center] at (1.25,-1.9) {\normalsize \textbf{às vezes inventa coisas!}};
        \node[align=center] at (1.25,-2.2) {\small Como a tia que conta fofoca};
        \node[align=center] at (1.25,-2.5) {\small sem ter certeza};
        \node[align=center] at (1.25,-2.8) {\tiny Pode parecer verdade, mas não é!};
    \end{tikzpicture}
    \end{center}
\end{frame}

\begin{frame}
    \frametitle{Por que "Tia Fofoqueira"?}
    \begin{center}
    \begin{tikzpicture}[scale=0.9]
        % Comparação 1
        \filldraw[fill=azul!20, draw=azul, very thick, rounded corners=12pt] (-5,2.2) rectangle (5,4);
        \node[align=center] at (0,3.5) {\normalsize \textbf{1. Sabe de TUDO}};
        \node[align=center] at (0,3.1) {\small (ou quase tudo)};
        \node[align=center] at (-2.5,2.5) {\small Tia: conhece toda família};
        \node[align=center] at (2.5,2.5) {\small ChatGPT: sabe de muitos assuntos};
        
        % Comparação 2
        \filldraw[fill=verde!20, draw=verde, very thick, rounded corners=12pt] (-5,0.5) rectangle (5,1.9);
        \node[align=center] at (0,1.4) {\normalsize \textbf{2. Adora Conversar}};
        \node[align=center] at (-2.5,0.8) {\small Tia: nunca cansa de falar};
        \node[align=center] at (2.5,0.8) {\small ChatGPT: responde sempre};
        
        % Comparação 3
        \filldraw[fill=laranja!20, draw=laranja, very thick, rounded corners=12pt] (-5,-1.2) rectangle (5,0.2);
        \node[align=center] at (0,-0.3) {\normalsize \textbf{3. Está Sempre Disponível}};
        \node[align=center] at (-2.5,-0.8) {\small Tia: no WhatsApp 24h};
        \node[align=center] at (2.5,-0.8) {\small ChatGPT: nunca dorme};
        
        % Comparação 4
        \filldraw[fill=amarelo!25, draw=orange, very thick, rounded corners=12pt] (-5,-3.2) rectangle (5,-1.3);
        \node[align=center] at (0,-1.9) {\normalsize \textbf{4. Mas Cuidado:}};
        \node[align=center] at (0,-2.3) {\normalsize \textbf{Pode Inventar!}};
        \node[align=center] at (-2.5,-2.8) {\small Tia: conta fofoca};
        \node[align=center] at (-2.5,-3.1) {\small sem certeza};
        \node[align=center] at (2.5,-2.8) {\small ChatGPT: pode};
        \node[align=center] at (2.5,-3.1) {\small estar errado};
    \end{tikzpicture}
    \end{center}
\end{frame}

\begin{frame}
    \frametitle{Exemplos: Tia Fofoqueira vs. ChatGPT}
    \begin{center}
    \begin{tikzpicture}[scale=0.85]
        % Exemplo 1
        \filldraw[fill=azul!20, draw=azul, very thick, rounded corners=12pt] (-5,2.3) rectangle (5,3.8);
        \node[align=center] at (0,3.4) {\normalsize \textbf{Exemplo 1: Sabe de Tudo}};
        \node[align=center] at (-2.5,2.8) {\small Tia: "Eu sei quem};
        \node[align=center] at (-2.5,2.5) {\small namora quem"};
        \node[align=center] at (2.5,2.8) {\small ChatGPT: "Eu sei sobre};
        \node[align=center] at (2.5,2.5) {\small matemática"};
        
        % Exemplo 2
        \filldraw[fill=verde!20, draw=verde, very thick, rounded corners=12pt] (-5,0.6) rectangle (5,2.1);
        \node[align=center] at (0,1.7) {\normalsize \textbf{Exemplo 2: Adora Conversar}};
        \node[align=center] at (-2.5,1.1) {\small Tia: "Me conta tudo!"};
        \node[align=center] at (2.5,1.3) {\small ChatGPT: "Pode perguntar};
        \node[align=center] at (2.5,0.9) {\small qualquer coisa!"};
        
        % Exemplo 3
        \filldraw[fill=laranja!20, draw=laranja, very thick, rounded corners=12pt] (-5,-1.1) rectangle (5,0.4);
        \node[align=center] at (0,-0.1) {\normalsize \textbf{Exemplo 3: Sempre Disponível}};
        \node[align=center] at (-2.5,-0.6) {\small Tia: "Me chama a};
        \node[align=center] at (-2.5,-0.9) {\small qualquer hora"};
        \node[align=center] at (2.5,-0.6) {\small ChatGPT: "Estou sempre};
        \node[align=center] at (2.5,-0.9) {\small online"};
        
        % Exemplo 4
        \filldraw[fill=amarelo!25, draw=orange, very thick, rounded corners=12pt] (-5,-2.8) rectangle (5,-1.4);
        \node[align=center] at (0,-1.9) {\normalsize \textbf{Exemplo 4: Pode Inventar}};
        \node[align=center] at (-2.5,-2.4) {\small Tia: "Ouvi dizer que..."};
        \node[align=center] at (2.5,-2.4) {\small ChatGPT: "Pode estar};
        \node[align=center] at (2.5,-2.7) {\small errado"};
    \end{tikzpicture}
    \end{center}
    
    \vspace{0.3cm}
    \begin{center}
        \normalsize \textbf{Por isso: sempre valide o que o ChatGPT diz!}
    \end{center}
\end{frame}

\begin{frame}
    \frametitle{O que é o ChatGPT?}
    \begin{center}
    \begin{tikzpicture}[scale=0.8]
        % Robô cabeça
        \filldraw[fill=azulclaro!30, draw=azul, very thick, rounded corners=10pt] (-2,1) rectangle (2,3);
        % Olhos
        \filldraw[fill=green, draw=darkgreen, very thick] (-0.8,2.2) circle (0.3cm);
        \filldraw[fill=green, draw=darkgreen, very thick] (0.8,2.2) circle (0.3cm);
        % Antena
        \draw[very thick, azul] (0,3) -- (0,3.5);
        \filldraw[fill=vermelho, draw=darkred] (0,3.5) circle (0.15cm);
        % Corpo
        \filldraw[fill=azulclaro!30, draw=azul, very thick, rounded corners=10pt] (-1.5,-1) rectangle (1.5,1);
        % Botões
        \filldraw[fill=laranja, draw=orange] (-0.5,0) circle (0.2cm);
        \filldraw[fill=laranja, draw=orange] (0.5,0) circle (0.2cm);
    \end{tikzpicture}
    \end{center}
    
    \vspace{0.5cm}
    \begin{center}
        \Large \textbf{Robô super inteligente}\\
        \Large \textbf{Conversa com você}\\
        \Large \textbf{Grátis e sempre disponível!}
    \end{center}
\end{frame}

\begin{frame}
    \frametitle{Por que usar?}
    \begin{center}
    \begin{tikzpicture}
        % Relógio 24/7
        \draw[very thick, azul] (0,0) circle (1.5cm);
        \draw[very thick, azul] (0,0) -- (0,1);
        \draw[very thick, azul] (0,0) -- (0.8,-0.8);
        \node at (0,-2.5) {\Large \textbf{24/7}};
        
        % Paciência
        \draw[very thick, verde] (4,0) circle (1.5cm);
        \node at (4,0) {\Huge \textbf{∞}};
        \node at (4,-2.5) {\Large \textbf{Paciência}};
        
        % Grátis
        \draw[very thick, laranja] (8,0) circle (1.5cm);
        \node at (8,0) {\Huge \textbf{\$}};
        \node at (8,-1.5) {\Large \textbf{0}};
        \node at (8,-2.5) {\Large \textbf{Grátis}};
    \end{tikzpicture}
    \end{center}
\end{frame}

\section{Como Fazer Perguntas}

\begin{frame}
    \frametitle{Pergunta Ruim vs. Pergunta Boa}
    \begin{columns}
        \begin{column}{0.5\textwidth}
            \begin{center}
            \begin{tikzpicture}[scale=0.8]
                \filldraw[fill=vermelho!20, draw=vermelho, very thick, rounded corners=12pt] (-2.5,-2.5) rectangle (2.5,2.5);
                \node at (0,1.5) {\Large \textbf{X}};
                \node[align=center] at (0,0.3) {\normalsize \textbf{Explique}\\
                \normalsize \textbf{frações}};
                \node[align=center] at (0,-1.2) {\small Muito vago!\\
                \small Não sabe o que você precisa};
            \end{tikzpicture}
            \end{center}
        \end{column}
        \begin{column}{0.5\textwidth}
            \begin{center}
            \begin{tikzpicture}[scale=0.8]
                \filldraw[fill=verde!20, draw=verde, very thick, rounded corners=12pt] (-2.5,-2.5) rectangle (2.5,2.5);
                \node at (0,1.5) {\Large \textbf{OK}};
                \node[align=center] at (0,0.5) {\normalsize \textbf{8ª série}\\
                \normalsize \textbf{Explicação simples}\\
                \normalsize \textbf{Exemplo: pizza}};
                \node[align=center] at (0,-1.2) {\small ChatGPT sabe\\
                \small exatamente o que você precisa};
            \end{tikzpicture}
            \end{center}
        \end{column}
    \end{columns}
    
    \vspace{0.5cm}
    \begin{center}
        \Large \textbf{Seja específico!}
    \end{center}
\end{frame}

\begin{frame}
    \frametitle{Template Perfeito}
    \begin{center}
    \begin{tikzpicture}[scale=1.0]
        % Caixa principal
        \filldraw[fill=azulclaro!20, draw=azul, very thick, rounded corners=15pt] (-4.5,-3.5) rectangle (4.5,3.5);
        
        % Elementos do template
        \filldraw[fill=azul!30, draw=azul, very thick, rounded corners=10pt] (-4,2) rectangle (4,2.8);
        \node[align=left] at (0,2.4) {\normalsize \textbf{1. Sua série:} 8ª série};
        
        \filldraw[fill=verde!30, draw=verde, very thick, rounded corners=10pt] (-4,0.8) rectangle (4,1.6);
        \node[align=left] at (0,1.2) {\normalsize \textbf{2. O problema:} Não entendi frações};
        
        \filldraw[fill=laranja!30, draw=laranja, very thick, rounded corners=10pt] (-4,-0.4) rectangle (4,0.4);
        \node[align=left] at (0,0) {\normalsize \textbf{3. Como explicar:} Como se tivesse 10 anos};
        
        \filldraw[fill=roxo!30, draw=roxo, very thick, rounded corners=10pt] (-4,-1.6) rectangle (4,-0.8);
        \node[align=left] at (0,-1.2) {\normalsize \textbf{4. Exemplo:} Use analogia da pizza};
    \end{tikzpicture}
    \end{center}
    
    \vspace{0.3cm}
    \begin{center}
        \Large \textbf{Quanto mais detalhes, melhor!}
    \end{center}
\end{frame}

\begin{frame}
    \frametitle{Dicas Visuais}
    \begin{center}
    \begin{tikzpicture}[scale=1.0]
        % Dica 1 - Específico
        \filldraw[fill=azul!30, draw=azul, very thick, rounded corners=12pt] (-3.5,2) rectangle (3.5,3.8);
        \node[align=center] at (0,3.1) {\Large \textbf{1. Seja Específico}};
        \node[align=center] at (0,2.4) {\small Diga exatamente o que você precisa};
        
        % Dica 2 - Contexto
        \filldraw[fill=verde!30, draw=verde, very thick, rounded corners=12pt] (-3.5,0) rectangle (3.5,1.8);
        \node[align=center] at (0,1.1) {\Large \textbf{2. Conte o Contexto}};
        \node[align=center] at (0,0.4) {\small Explique a situação};
        
        % Dica 3 - Simples
        \filldraw[fill=laranja!30, draw=laranja, very thick, rounded corners=12pt] (-3.5,-2) rectangle (3.5,-0.2);
        \node[align=center] at (0,-0.9) {\Large \textbf{3. Peça Explicação Simples}};
        \node[align=center] at (0,-1.4) {\small Como se você fosse mais novo};
        
        % Dica 4 - Exemplos
        \filldraw[fill=roxo!30, draw=roxo, very thick, rounded corners=12pt] (-3.5,-4) rectangle (3.5,-2.2);
        \node[align=center] at (0,-2.9) {\Large \textbf{4. Peça Exemplos Práticos}};
        \node[align=center] at (0,-3.4) {\small Analogias do dia a dia};
    \end{tikzpicture}
    \end{center}
\end{frame}

\section{Como Usar para Estudar}

\begin{frame}
    \frametitle{Quando Usar?}
    \begin{center}
    \begin{tikzpicture}[scale=0.9]
        % Não entendeu
        \filldraw[fill=vermelho!20, draw=vermelho, very thick, rounded corners=12pt] (-4.5,1) rectangle (-0.5,3.5);
        \node at (-2.5,2.8) {\Large \textbf{?}};
        \node[align=center] at (-2.5,1.8) {\normalsize \textbf{Não}\\
        \normalsize \textbf{entendeu}};
        
        % Exercício
        \filldraw[fill=laranja!20, draw=laranja, very thick, rounded corners=12pt] (-0.2,1) rectangle (3.2,3.5);
        \node at (1.5,2.8) {\Large \textbf{E}};
        \node[align=center] at (1.5,1.8) {\normalsize \textbf{Exercício}\\
        \small travado};
        
        % Revisão
        \filldraw[fill=verde!20, draw=verde, very thick, rounded corners=12pt] (3.5,1) rectangle (7,3.5);
        \node at (5.25,2.8) {\Large \textbf{R}};
        \node[align=center] at (5.25,1.8) {\normalsize \textbf{Revisar}\\
        \small para prova};
        
        % Setas para ChatGPT
        \draw[->, very thick, azul] (-2.5,0.5) -- (-2.5,-0.5);
        \draw[->, very thick, azul] (1.5,0.5) -- (1.5,-0.5);
        \draw[->, very thick, azul] (5.25,0.5) -- (5.25,-0.5);
        
        % ChatGPT
        \filldraw[fill=azul!30, draw=azul, very thick, rounded corners=15pt] (-3.5,-2.5) rectangle (7.5,-0.5);
        \node at (2,-1.5) {\Large \textbf{ChatGPT}};
    \end{tikzpicture}
    \end{center}
\end{frame}

\begin{frame}
    \frametitle{Exemplos Práticos}
    \begin{center}
    \begin{tikzpicture}[scale=0.85]
        % Exemplo 1
        \filldraw[fill=azul!20, draw=azul, very thick, rounded corners=12pt] (-4.5,2) rectangle (4.5,4.2);
        \node[align=center] at (0,3.5) {\normalsize \textbf{Exemplo 1: Não Entendeu}};
        \node[align=center] at (0,2.6) {\small "Explica fotossíntese como se eu tivesse 10 anos"};
        
        % Exemplo 2
        \filldraw[fill=verde!20, draw=verde, very thick, rounded corners=12pt] (-4.5,-0.3) rectangle (4.5,1.9);
        \node[align=center] at (0,1.2) {\normalsize \textbf{Exemplo 2: Exercício}};
        \node[align=center] at (0,0.3) {\small "Me dá dicas, não a resposta!"};
        
        % Exemplo 3
        \filldraw[fill=laranja!20, draw=laranja, very thick, rounded corners=12pt] (-4.5,-2.8) rectangle (4.5,-0.6);
        \node[align=center] at (0,-1.5) {\normalsize \textbf{Exemplo 3: Revisão}};
        \node[align=center] at (0,-2.2) {\small "Faz resumo curto para prova"};
    \end{tikzpicture}
    \end{center}
\end{frame}

\section{O que Fazer e NÃO Fazer}

\begin{frame}
    \frametitle{FAÇA vs. NÃO FAÇA}
    \begin{center}
    \begin{tikzpicture}[scale=1.0]
        % Lado FAÇA
        \filldraw[fill=verde!30, draw=verde, very thick, rounded corners=15pt] (-5.5,-3.5) rectangle (-0.5,3.5);
        \node at (-3,2.5) {\Large \textbf{OK FAÇA}};
        \node[align=left] at (-3,0.8) {\normalsize • Seja específico};
        \node[align=left] at (-3,0) {\normalsize • Peça explicações};
        \node[align=left] at (-3,-0.8) {\normalsize • Use para aprender};
        \node[align=left] at (-3,-1.6) {\normalsize • Valide respostas};
        
        % Seta
        \draw[<->, very thick, azul] (0,0) -- (0.5,0);
        
        % Lado NÃO FAÇA
        \filldraw[fill=vermelho!30, draw=vermelho, very thick, rounded corners=15pt] (1,-3.5) rectangle (6,3.5);
        \node at (3.5,2.5) {\Large \textbf{X NÃO FAÇA}};
        \node[align=left] at (3.5,0.8) {\normalsize • Copiar sem entender};
        \node[align=left] at (3.5,0) {\normalsize • Confiar cegamente};
        \node[align=left] at (3.5,-0.8) {\normalsize • Usar para trapacear};
        \node[align=left] at (3.5,-1.6) {\normalsize • Pular estudo};
    \end{tikzpicture}
    \end{center}
\end{frame}

\begin{frame}
    \frametitle{Regra de Ouro}
    \begin{center}
    \begin{tikzpicture}[scale=1.2]
        % Círculo dourado
        \filldraw[fill=yellow!40, draw=orange, very thick] (0,0) circle (2.5cm);
        \node at (0,0.5) {\Huge \textbf{APRENDER}};
        \node at (0,-0.5) {\Huge \textbf{NÃO}};
        \node at (0,-1.3) {\Large \textbf{EVITAR}};
        
        % Setas
        \draw[->, very thick, verde] (-3.5,0) -- (-2.8,0);
        \draw[->, very thick, vermelho] (2.8,0) -- (3.5,0);
        
        \node at (-4.5,0) {\Large \textbf{OK}};
        \node at (4.5,0) {\Large \textbf{X}};
    \end{tikzpicture}
    \end{center}
    
    \vspace{0.5cm}
    \begin{center}
        \Large \textbf{Use para aprender, não para evitar aprender!}
    \end{center}
\end{frame}

\section{Cuidados}

\begin{frame}
    \frametitle{Cuidado! Pode estar errado}
    \begin{center}
    \begin{tikzpicture}[scale=1.0]
        % Cabeça com interrogação
        \filldraw[fill=amarelo!30, draw=orange, very thick] (0,1.5) circle (1.8cm);
        \node at (0,1.5) {\Huge \textbf{?}};
        
        % Problemas
        \filldraw[fill=vermelho!20, draw=vermelho, very thick, rounded corners=12pt] (-4,-1.2) rectangle (-1.5,-0.2);
        \node[align=center] at (-2.75,-0.7) {\normalsize \textbf{Pode inventar}\\
        \small informações};
        
        \filldraw[fill=laranja!20, draw=laranja, very thick, rounded corners=12pt] (-0.7,-1.2) rectangle (1.7,-0.2);
        \node[align=center] at (0.5,-0.7) {\normalsize \textbf{Pode estar}\\
        \small desatualizado};
        
        \filldraw[fill=azul!20, draw=azul, very thick, rounded corners=12pt] (2.5,-1.2) rectangle (5,-0.2);
        \node[align=center] at (3.75,-0.7) {\normalsize \textbf{Pode confundir}\\
        \small conceitos};
        
        % Solução
        \filldraw[fill=verde!30, draw=verde, very thick, rounded corners=12pt] (-2.5,-3.2) rectangle (2.5,-2.2);
        \node at (0,-2.7) {\Large \textbf{Sempre valide!}};
    \end{tikzpicture}
    \end{center}
\end{frame}

\begin{frame}
    \frametitle{Quando NÃO usar}
    \begin{center}
    \begin{tikzpicture}[scale=0.8]
        % Proibido em provas
        \draw[very thick, vermelho] (-3,2) circle (1cm);
        \draw[very thick, vermelho] (-3.7,1.3) -- (-2.3,2.7);
        \node at (-3,-0.5) {\small \textbf{Provas}};
        
        % Proibido copiar
        \draw[very thick, vermelho] (0,2) circle (1cm);
        \draw[very thick, vermelho] (-0.7,1.3) -- (0.7,2.7);
        \node at (0,-0.5) {\small \textbf{Copiar}};
        
        % Proibido substituir
        \draw[very thick, vermelho] (3,2) circle (1cm);
        \draw[very thick, vermelho] (2.3,1.3) -- (3.7,2.7);
        \node at (3,-0.5) {\small \textbf{Substituir}};
        
        % Quando usar
        \filldraw[fill=verde!30, draw=verde, very thick, rounded corners=10pt] (-3.5,-2.5) rectangle (3.5,-1.5);
        \node at (0,-2) {\Large \textbf{Use para: Entender, Praticar, Revisar}};
    \end{tikzpicture}
    \end{center}
\end{frame}

\section{Exemplos}

\begin{frame}
    \frametitle{Exemplo 1: Não Entendeu}
    \begin{center}
    \begin{tikzpicture}[scale=0.9]
        % Problema
        \filldraw[fill=vermelho!20, draw=vermelho, very thick, rounded corners=12pt] (-4.5,2) rectangle (4.5,3.8);
        \node[align=center] at (0,3.1) {\normalsize \textbf{Problema:} Não entendi fotossíntese};
        
        % Pergunta
        \filldraw[fill=azul!20, draw=azul, very thick, rounded corners=12pt] (-4.5,0) rectangle (4.5,1.8);
        \node[align=center] at (0,1.1) {\small "Explica fotossíntese como se eu tivesse 10 anos,\\
        usando analogia de fábrica"};
        
        % Solução
        \filldraw[fill=verde!20, draw=verde, very thick, rounded corners=12pt] (-4.5,-2.2) rectangle (4.5,-0.4);
        \node[align=center] at (0,-1.4) {\normalsize \textbf{Resultado:} Entendeu de verdade!};
        
        % Setas
        \draw[->, very thick, azul] (0,1.8) -- (0,2);
        \draw[->, very thick, verde] (0,0) -- (0,-0.2);
    \end{tikzpicture}
    \end{center}
\end{frame}

\begin{frame}
    \frametitle{Exemplo 2: Exercício}
    \begin{center}
    \begin{tikzpicture}[scale=0.9]
        % Problema
        \filldraw[fill=laranja!20, draw=laranja, very thick, rounded corners=12pt] (-4.5,2) rectangle (4.5,3.8);
        \node[align=center] at (0,3.1) {\normalsize \textbf{Problema:} Travado no exercício};
        
        % Pergunta
        \filldraw[fill=azul!20, draw=azul, very thick, rounded corners=12pt] (-4.5,0) rectangle (4.5,1.8);
        \node[align=center] at (0,1.1) {\small "Me dá dicas passo a passo,\\
        mas não me dê a resposta!"};
        
        % Solução
        \filldraw[fill=verde!20, draw=verde, very thick, rounded corners=12pt] (-4.5,-2.2) rectangle (4.5,-0.4);
        \node[align=center] at (0,-1.4) {\normalsize \textbf{Resultado:} Aprendeu o método!};
        
        % Setas
        \draw[->, very thick, azul] (0,1.8) -- (0,2);
        \draw[->, very thick, verde] (0,0) -- (0,-0.2);
    \end{tikzpicture}
    \end{center}
\end{frame}

\begin{frame}
    \frametitle{Exemplo 3: Revisão}
    \begin{center}
    \begin{tikzpicture}[scale=0.9]
        % Problema
        \filldraw[fill=roxo!20, draw=roxo, very thick, rounded corners=12pt] (-4.5,2) rectangle (4.5,3.8);
        \node[align=center] at (0,3.1) {\normalsize \textbf{Problema:} Prova amanhã, muita matéria};
        
        % Pergunta
        \filldraw[fill=azul!20, draw=azul, very thick, rounded corners=12pt] (-4.5,0) rectangle (4.5,1.8);
        \node[align=center] at (0,1.1) {\small "Faz resumo super curto\\
        só o essencial"};
        
        % Solução
        \filldraw[fill=verde!20, draw=verde, very thick, rounded corners=12pt] (-4.5,-2.2) rectangle (4.5,-0.4);
        \node[align=center] at (0,-1.4) {\normalsize \textbf{Resultado:} Revisou rápido!};
        
        % Setas
        \draw[->, very thick, azul] (0,1.8) -- (0,2);
        \draw[->, very thick, verde] (0,0) -- (0,-0.2);
    \end{tikzpicture}
    \end{center}
\end{frame}

\section{Conclusão}

\begin{frame}
    \frametitle{Resumo Visual}
    \begin{center}
    \begin{tikzpicture}[scale=0.7]
        % Tia fofoqueira
        \filldraw[fill=laranja!30, draw=orange, very thick] (-4,2) circle (0.8cm);
        \node at (-4,2) {\Large \textbf{C}};
        \node at (-4,0.8) {\small Tia};
        \node at (-4,0.3) {\small Fofoqueira};
        
        % Perguntas boas
        \filldraw[fill=verde!30, draw=verde, very thick, rounded corners=10pt] (-1,1) rectangle (1,3);
        \node at (0,2.5) {\Large \textbf{OK}};
        \node at (0,1.5) {\small Perguntas};
        \node at (0,1.2) {\small Boas};
        
        % Aprender
        \filldraw[fill=azul!30, draw=azul, very thick, rounded corners=10pt] (2.5,1) rectangle (4.5,3);
        \node at (3.5,2.5) {\Large \textbf{L}};
        \node at (3.5,1.5) {\small Aprender};
        
        % Validar
        \filldraw[fill=amarelo!30, draw=orange, very thick, rounded corners=10pt] (-1,-1.5) rectangle (1,0.5);
        \node at (0,0) {\Large \textbf{OK}};
        \node at (0,-0.8) {\small Validar};
        
        % Não trapacear
        \filldraw[fill=vermelho!30, draw=vermelho, very thick, rounded corners=10pt] (2.5,-1.5) rectangle (4.5,0.5);
        \node at (3.5,0) {\Large \textbf{X}};
        \node at (3.5,-0.8) {\small Trapacear};
    \end{tikzpicture}
    \end{center}
\end{frame}

\begin{frame}
    \frametitle{Dicas Finais}
    \begin{center}
    \begin{tikzpicture}[scale=1.0]
        % Dica 1
        \filldraw[fill=azul!30, draw=azul, very thick, rounded corners=12pt] (-4,2) rectangle (-1,3.8);
        \node[align=center] at (-2.5,3.2) {\normalsize \textbf{1.}\\
        \small Experimente};
        
        % Dica 2
        \filldraw[fill=verde!30, draw=verde, very thick, rounded corners=12pt] (-0.3,2) rectangle (2.3,3.8);
        \node[align=center] at (1,3.2) {\normalsize \textbf{2.}\\
        \small Valide};
        
        % Dica 3
        \filldraw[fill=laranja!30, draw=laranja, very thick, rounded corners=12pt] (2.7,2) rectangle (5,3.8);
        \node[align=center] at (3.85,3.2) {\normalsize \textbf{3.}\\
        \small Pratique};
        
        % Mensagem central
        \filldraw[fill=amarelo!40, draw=orange, very thick, rounded corners=15pt] (-4,-1.2) rectangle (5,1.2);
        \node[align=center] at (0.5,0) {\Large \textbf{Use com sabedoria!}\\
        \normalsize É seu ajudante, não substituto!};
    \end{tikzpicture}
    \end{center}
\end{frame}

\begin{frame}
    \frametitle{Dúvidas?}
    \begin{center}
    \begin{tikzpicture}[scale=1.2]
        % Interrogação grande
        \filldraw[fill=azul!30, draw=azul, very thick] (0,0) circle (2cm);
        \node at (0,0) {\Huge \textbf{?}};
    \end{tikzpicture}
    \end{center}
    
    \vspace{0.5cm}
    \begin{center}
        \Large \textbf{Obrigado pela atenção!}\\
        \vspace{0.3cm}
        \normalsize \textbf{Use o ChatGPT com sabedoria!}
    \end{center}
    
    \vspace{0.3cm}
    \begin{center}
        \small Professor Dr. Bento Rafael Siqueira\\
        \small bento.siqueira@ufla.br
    \end{center}
\end{frame}

\end{document}